\documentclass[UTF8]{ctexart} 
\usepackage[a4paper,margin=2.6cm]{geometry}
\usepackage{amsmath, amssymb, amsthm}
\usepackage{graphicx}
\usepackage[colorlinks=true,linkcolor=blue]{hyperref}

\title{复变函数学习笔记：3.4节“幅角原理及其应用”}
\date{}

\newtheorem{theorem}{定理}[section]
\newtheorem{lemma}[theorem]{引理}
\newtheorem{definition}[theorem]{定义}
\newtheorem{corollary}[theorem]{推论}
\newtheorem{remark}{注记}[section]
\numberwithin{equation}{section}

\newcommand{\RR}{\mathbb{R}}
\newcommand{\CC}{\mathbb{C}}

\begin{document}

\maketitle

\section{引言}

在复分析中，一个自然的问题是如何计算一个解析函数在给定区域内零点的个数（计重数），以及极点个数（计阶数）。Cauchy 留数定理为我们提供了一种有力的工具：若我们能找到某个函数的积分，其留数恰好反映零点和极点的信息，那么问题便转化为计算围道积分。观察到若 $f$ 在 $z_0$ 处有 $n$ 阶零点，则 $f(z)=(z-z_0)^n g(z)$，$g(z_0)\neq0$，从而
\[
\frac{f'(z)}{f(z)} = \frac{n}{z-z_0} + \frac{g'(z)}{g(z)},
\]
即 $f'/f$ 在零点处有一阶极点，留数为 $n$。类似地，若 $z_0$ 是 $m$ 阶极点，则 $f'/f$ 的留数为 $-m$。因此，函数 $f'/f$（称为 $f$ 的\textbf{对数导数}）的围道积分将直接给出内部零点与极点的代数和。这一思想正是\textbf{幅角原理}的核心。

\section{对数导数与幅角原理}

\begin{definition}[对数导数]
设 $f$ 在区域 $\Omega$ 内亚纯，且 $f$ 不恒为零。则函数
\[
\frac{f'(z)}{f(z)}
\]
称为 $f$ 的\textbf{对数导数}。其名称来源于 $( \log f(z) )' = f'(z)/f(z)$，但 $\log f$ 可能是多值的，而 $f'/f$ 是单值解析函数（除去 $f$ 的零点和极点）。
\end{definition}

\begin{theorem}[幅角原理]
设 $f$ 在开集 $\Omega$ 内亚纯，$\gamma$ 是 $\Omega$ 内一条可求长的简单闭曲线，其内部 $\operatorname{Int}(\gamma)$ 完全包含在 $\Omega$ 中。假设 $f$ 在 $\gamma$ 上既无零点也无极点。则
\[
\frac{1}{2\pi i}\oint_{\gamma} \frac{f'(z)}{f(z)}\,dz = N - P,
\]
其中 $N$ 是 $f$ 在 $\gamma$ 内部零点的个数（计重数），$P$ 是 $f$ 在 $\gamma$ 内部极点的个数（计阶数）。
\end{theorem}

\begin{proof}
由留数定理（Corollary 2.3, Chapter 3），左边等于 $\gamma$ 内部所有奇点的留数之和。$f'/f$ 的奇点恰好是 $f$ 的零点和极点。若 $z_0$ 是 $f$ 的 $n$ 阶零点，则在 $z_0$ 附近 $f(z)=(z-z_0)^n g(z)$，$g$ 解析且 $g(z_0)\neq0$。于是
\[
\frac{f'(z)}{f(z)} = \frac{n}{z-z_0} + \frac{g'(z)}{g(z)},
\]
右边第二项在 $z_0$ 解析，故 $\operatorname{Res}_{z_0}(f'/f)=n$。若 $z_0$ 是 $f$ 的 $m$ 阶极点，则 $f(z)=(z-z_0)^{-m}h(z)$，$h$ 解析且 $h(z_0)\neq0$，类似可得 $\operatorname{Res}_{z_0}(f'/f)=-m$。因此所有奇点的留数和为 $N-P$，定理得证。
\end{proof}

\begin{remark}[几何解释]
注意到 $\frac{f'(z)}{f(z)}dz = d(\log f(z)) = d(\ln|f(z)| + i\arg f(z))$。由于 $|f(z)|$ 是单值的，沿闭曲线 $\gamma$ 积分后其实部变化为零，而虚部给出 $\arg f(z)$ 的变化量。因此
\[
\frac{1}{2\pi i}\oint_{\gamma} \frac{f'(z)}{f(z)}\,dz = \frac{1}{2\pi} \Delta_{\gamma} \arg f(z).
\]
即积分值等于当 $z$ 沿 $\gamma$ 绕行一周时，$f(z)$ 的辐角改变量除以 $2\pi$，也就是 $f(\gamma)$ 绕原点的圈数（ winding number）。幅角原理断言：这个圈数恰好等于内部零点数减去极点数。
\end{remark}

\section{Rouché 定理}

幅角原理的一个直接而有力的推论是 Rouché 定理，它给出了判断两个函数在区域内零点个数相等的简单条件。

\begin{theorem}[Rouché 定理]
设 $f$ 和 $g$ 在区域 $\Omega$ 内解析，$\gamma$ 是 $\Omega$ 内一条可求长的简单闭曲线，其内部 $\operatorname{Int}(\gamma) \subset \Omega$。若在 $\gamma$ 上
\[
|f(z)| > |g(z)|,
\]
则 $f$ 与 $f+g$ 在 $\gamma$ 内部有相同个数的零点（计重数）。
\end{theorem}

\subsection{证明思路}

Rouché 定理的核心思想是：当我们在 $f$ 上添加一个在边界上“足够小”的扰动 $g$ 时，$f+g$ 的零点个数不会改变。这可以通过考虑函数族 $F_t(z)=f(z)+t g(z)$，$t\in[0,1]$ 来实现。条件 $|f|>|g|$ 保证了 $F_t$ 在 $\gamma$ 上无零点，因此 $F_t$ 的辐角沿 $\gamma$ 连续变化，其绕原点的圈数 $N(t)$ 是 $t$ 的连续整数函数，从而为常数。于是 $N(0)=N(1)$，即 $f$ 与 $f+g$ 有相同的零点个数。

\subsection{详细证明}

\begin{proof}
对于 $t\in[0,1]$，定义 $F_t(z)=f(z)+t g(z)$。在 $\gamma$ 上，由假设 $|f(z)| > |g(z)| \ge |t g(z)|$，故 $|F_t(z)| \ge |f(z)| - |t g(z)| > 0$，因此 $F_t$ 在 $\gamma$ 上无零点。令
\[
N(t) = \frac{1}{2\pi i}\oint_{\gamma} \frac{F_t'(z)}{F_t(z)}\,dz.
\]
由幅角原理，$N(t)$ 是整数，等于 $F_t$ 在 $\gamma$ 内部零点的个数（因为 $F_t$ 解析，无极点）。被积函数关于 $(z,t)$ 连续，且 $F_t(z)$ 在 $\gamma$ 上一致远离零，故 $N(t)$ 是 $t$ 的连续函数。连续整数函数必为常数，所以 $N(0)=N(1)$，即 $f$ 与 $f+g$ 在 $\gamma$ 内部零点个数相同。
\end{proof}

\begin{remark}
Rouché 定理常被用来估计多项式根的分布。例如，对于多项式 $P(z)=a_n z^n + \cdots + a_0$，取 $f(z)=a_n z^n$，$g(z)=a_{n-1}z^{n-1}+\cdots+a_0$，若在某个大圆 $|z|=R$ 上 $|a_n|R^n > |a_{n-1}|R^{n-1}+\cdots+|a_0|$，则 $P$ 在该圆内有 $n$ 个根（计重数），这正是代数学基本定理的另一种证明。
\end{remark}

\subsection{思想来源}

Rouché 定理的思想可追溯到 Cauchy 的幅角原理以及“扰动不改变指标”的直观。Eugène Rouché 在 1862 年给出了这一定理。它本质上是幅角原理的一个同伦版本：函数族 $F_t$ 给出了从 $f$ 到 $f+g$ 的一个连续形变，只要在形变过程中 $F_t(\gamma)$ 始终不穿过原点，那么 $F_t$ 绕原点的圈数就保持不变。

\section{应用举例}

\subsection{例1：多项式根的位置}
判断 $P(z)=z^5+3z+1$ 在单位圆 $|z|=1$ 内根的个数。取 $f(z)=3z$，$g(z)=z^5+1$。在 $|z|=1$ 上，$|f(z)|=3$，$|g(z)|\le |z|^5+1=2$，故 $|f|>|g|$。$f$ 在单位圆内只有一个零点 $z=0$，因此 $P$ 在单位圆内恰有一个根。

\subsection{例2：辐角原理求零点个数}
设 $f(z)=z^4+2z^2+1$，在圆 $|z|=2$ 上，$|f(z)-z^4| = |2z^2+1| \le 2\cdot 4+1=9$，而 $|z^4|=16$，故 $|f|>|f-z^4|$ 不成立（因为 $9<16$），不能直接用 Rouché。但我们可以直接计算 $f(z)=(z^2+1)^2$，有四个重根 $z=\pm i$。验证幅角原理：沿 $|z|=2$ 计算 $f'/f$ 的积分，结果应为 $4$，与零点个数一致。


在前面的部分中，我们已经介绍了幅角原理（Argument Principle）以及它的两个重要推论：鲁歇定理（Rouché's Theorem）和开映射定理（Open Mapping Theorem）。本节我们将继续讨论幅角原理的另一个深刻应用——最大模原理（Maximum Modulus Principle），并通过反例说明其成立的条件。

    % 让下一个定理编号为3.4.5
\begin{theorem}[最大模原理]
设 $\Omega \subset \CC$ 是一个区域，$f$ 是 $\Omega$ 上的非常值全纯函数。则 $|f(z)|$ 不可能在 $\Omega$ 的内部取得最大值。
\end{theorem}

\begin{proof}[证明思路]
该定理的证明非常简洁，其核心思想是利用\textbf{开映射定理}：非常值全纯函数将开集映为开集。假设 $|f|$ 在某点 $z_0 \in \Omega$ 取得最大值，那么 $f$ 在 $z_0$ 附近的像将包含一个以 $f(z_0)$ 为中心的小开圆盘。在该圆盘中必然存在模长大于 $|f(z_0)|$ 的点，这与最大值假设矛盾。
\end{proof}

\begin{proof}[详细证明]
假设存在 $z_0 \in \Omega$ 使得
\[
|f(z_0)| = \max_{z \in \Omega} |f(z)|.
\]
因为 $f$ 非常值且全纯，根据开映射定理（定理3.4.4），$f$ 将开集映为开集。特别地，对于任意充分小的圆盘 $D \subset \Omega$ 且 $z_0 \in D$，其像 $f(D)$ 是包含 $f(z_0)$ 的一个开集。因此，存在某个 $w \in f(D)$ 使得 $|w| > |f(z_0)|$。但 $w = f(z)$ 对于某个 $z \in D$，这就意味着 $|f(z)| > |f(z_0)|$，与 $z_0$ 是最大值点矛盾。
\end{proof}

最大模原理指出，非常值全纯函数的模长在区域内部不能达到极大值。结合连续性，我们立即得到关于闭包上最大值位置的结论：

\begin{corollary}
设 $\Omega$ 是一个有界区域，且其闭包 $\overline{\Omega}$ 是紧集。若 $f$ 在 $\Omega$ 上全纯，在 $\overline{\Omega}$ 上连续，则 $|f|$ 的最大值必定在边界 $\partial \Omega = \overline{\Omega} \setminus \Omega$ 上取得。即
\[
\sup_{z \in \Omega} |f(z)| \le \sup_{z \in \partial \Omega} |f(z)|.
\]
\end{corollary}

\begin{proof}
由于 $|f|$ 在紧集 $\overline{\Omega}$ 上连续，它必然在某点 $z_0 \in \overline{\Omega}$ 达到最大值。若 $f$ 是常数，结论平凡成立。若 $f$ 非常值，由最大模原理，$z_0$ 不可能在 $\Omega$ 内部，故 $z_0 \in \partial \Omega$。
\end{proof}

\begin{remark}[无界区域的反例]
推论中要求 $\overline{\Omega}$ 是紧集（即 $\Omega$ 有界）是必不可少的。对于无界区域，即使函数在边界上有界，内部也可能无界。一个经典的例子是：取第一象限 $\Omega = \{ z = x+iy : x>0, y>0 \}$，考虑函数 $F(z) = e^{-iz^2}$。$F$ 是整函数，在 $\overline{\Omega}$ 上连续。在边界正实轴和正虚轴上，分别有 $z=x$ 时 $|F(x)| = |e^{-ix^2}| = 1$，$z=iy$ 时 $|F(iy)| = |e^{iy^2}| = 1$。然而，若沿射线 $z = r e^{i\pi/4}$（$r>0$）取值，有 $F(z) = e^{-i(r e^{i\pi/4})^2} = e^{-i r^2 i} = e^{r^2}$，模长随 $r \to \infty$ 而趋于无穷。这说明 $|F|$ 在 $\Omega$ 内部无界，最大值不在边界上取得。
\end{remark}

最大模原理及其推论在复分析中有着广泛的应用，例如用于证明施瓦茨引理、刘维尔定理的变体以及单位圆盘上全纯函数的若干估计。至此，我们完成了第3.4节“幅角原理及其应用”的全部主要内容。




\end{document}